% GENERATED FILE. Edit canonical lessons, then run npm run generate:books.
% canonical-source-hash: fnv1a64:85b0fb183ae25112
% canonical-lessons: GU-C156-gharenu, GU-C156-callo, GU-C156-medi, GU-C156-sajavvu, GU-C156-badhvu

\chapter{Jewellery, bindi and henna}
\label{ch:gu-jewellery-bindi-and-henna}
\hlchaptermodality{156}

\begin{chapteropening}
\textbf{By the end of this chapter:} \emph{I can talk about jewellery, the bindi and henna.}

Complete the last lesson of chapter 156: I can talk about jewellery, the bindi and henna.
\end{chapteropening}

\section[\texorpdfstring{\gu{ઘરેણું} (ghareṇũ)}{ઘરેણું (ghareṇũ)}]{\gu{ઘરેણું} (ghareṇũ) — ornament, piece of jewellery}

\label{lesson:GU-C156-gharenu}



\begin{quote}
Before the new one: say the Gujarati for perfume, then the Gujarati for to apply.
\end{quote}

\subsection*{You'll want to know: \gu{ઘરેણું}}
\textbf{\gu{ઘરેણું}} — \emph{ghareṇũ} — \textquotedblleft{}ornament, piece of jewellery\textquotedblright{}.

\begin{quote}
\textbf{\gu{દાદી} \gu{પાસે} \gu{જૂનું} \gu{ઘરેણું} \gu{છે}.} — \emph{dādī pāse jūnũ ghareṇũ che} — \textquotedblleft{}Grandmother has an old piece of jewellery.\textquotedblright{}
\end{quote}

\begin{cousinweb}
\textbf{\gu{ઘરેણું}} (\emph{ghareṇũ}) is a tadbhava. CDIAL lists it, glossed \textquotedblleft{}jewelry\textquotedblright{}, under Sanskrit \emph{grahaṇa}, \textquotedblleft{}seizing, holding\textquotedblright{}, in its second sense, \textquotedblleft{}putting on (clothes)\textquotedblright{}; the same entry gives a Gujarati \emph{ghareṇe}, \textquotedblleft{}in pawn\textquotedblright{}, under the first sense (CDIAL 4364).
\end{cousinweb}

\begin{grammarlens}[title={What you've built}]
One more word for everyday life at A2.
\end{grammarlens}

\subsection*{Your turn}
\begin{itemize}
\raggedright
  \item \emph{Say it:} \emph{ghareṇũ}
  \item \emph{Say it:} \emph{ghareṇũ}, once more
  \item \emph{Say it:} \emph{dādī pāse jūnũ ghareṇũ che} — Grandmother has an old piece of jewellery.
  \item \emph{Recall:} say the Gujarati for perfume, then the Gujarati for to apply, then say \emph{ghareṇũ} again
\end{itemize}

\begin{tcolorbox}[breakable,colback=teal!4,colframe=teal!35!black,title={Before you move on}]
What does \gu{ઘરેણું} mean? (\textquotedblleft{}Ornament, piece of jewellery\textquotedblright{}.) Say it once more.
\end{tcolorbox}
\section[\texorpdfstring{\gu{ચાંલ્લો} (cā̃llo)}{ચાંલ્લો (cā̃llo)}]{\gu{ચાંલ્લો} (cā̃llo) — bindi, forehead mark}

\label{lesson:GU-C156-callo}



\begin{quote}
Before the new one: say the Gujarati for to apply, then the Gujarati for ornament.
\end{quote}

\subsection*{You'll want to know: \gu{ચાંલ્લો}}
\textbf{\gu{ચાંલ્લો}} — \emph{cā̃llo} — \textquotedblleft{}bindi, forehead mark\textquotedblright{}.

\begin{quote}
\textbf{\gu{તે} \gu{કપાળે} \gu{ચાંલ્લો} \gu{કરે} \gu{છે}.} — \emph{te kapāḷe cā̃llo kare che} — \textquotedblleft{}She puts a bindi on her forehead.\textquotedblright{}
\end{quote}

\begin{cousinweb}
The origin of \textbf{\gu{ચાંલ્લો}} (\emph{cā̃llo}) is not settled. CDIAL lists no Gujarati form of it, and EtymDB has no entry for it.
\end{cousinweb}

\begin{grammarlens}[title={What you've built}]
One more word for everyday life at A2.
\end{grammarlens}

\subsection*{Your turn}
\begin{itemize}
\raggedright
  \item \emph{Say it:} \emph{cā̃llo}
  \item \emph{Say it:} \emph{cā̃llo}, once more
  \item \emph{Say it:} \emph{te kapāḷe cā̃llo kare che} — She puts a bindi on her forehead.
  \item \emph{Recall:} say the Gujarati for to apply, then the Gujarati for ornament, then say \emph{cā̃llo} again
\end{itemize}

\begin{tcolorbox}[breakable,colback=teal!4,colframe=teal!35!black,title={Before you move on}]
What does \gu{ચાંલ્લો} mean? (\textquotedblleft{}Bindi, forehead mark\textquotedblright{}.) Say it once more.
\end{tcolorbox}
\section[\texorpdfstring{\gu{મેંદી} (mẽdī)}{મેંદી (mẽdī)}]{\gu{મેંદી} (mẽdī) — henna}

\label{lesson:GU-C156-medi}



\begin{quote}
Before the new one: say the Gujarati for ornament, then the Gujarati for bindi.
\end{quote}

\subsection*{You'll want to know: \gu{મેંદી}}
\textbf{\gu{મેંદી}} — \emph{mẽdī} — \textquotedblleft{}henna\textquotedblright{}.

\begin{quote}
\textbf{\gu{તહેવારમાં} \gu{છોકરીઓ} \gu{મેંદી} \gu{લગાડે} \gu{છે}.} — \emph{tahevārmā̃ chokrīo mẽdī lagāḍe che} — \textquotedblleft{}Girls put on henna at festivals.\textquotedblright{}
\end{quote}

\begin{cousinweb}
\textbf{\gu{મેંદી}} (\emph{mẽdī}) is a tadbhava. It comes from Sanskrit \emph{mēndhī}, \emph{mēndhikā}, the plant \emph{Lawsonia alba}, which gives henna, and CDIAL lists the Gujarati word with that sense (CDIAL 10328).
\end{cousinweb}

\begin{grammarlens}[title={What you've built}]
One more word for everyday life at A2.
\end{grammarlens}

\subsection*{Your turn}
\begin{itemize}
\raggedright
  \item \emph{Say it:} \emph{mẽdī}
  \item \emph{Say it:} \emph{mẽdī}, once more
  \item \emph{Say it:} \emph{tahevārmā̃ chokrīo mẽdī lagāḍe che} — Girls put on henna at festivals.
  \item \emph{Recall:} say the Gujarati for ornament, then the Gujarati for bindi, then say \emph{mẽdī} again
\end{itemize}

\begin{tcolorbox}[breakable,colback=teal!4,colframe=teal!35!black,title={Before you move on}]
What does \gu{મેંદી} mean? (\textquotedblleft{}Henna\textquotedblright{}.) Say it once more.
\end{tcolorbox}
\section[\texorpdfstring{\gu{સજાવવું} (sajāvvũ)}{સજાવવું (sajāvvũ)}]{\gu{સજાવવું} (sajāvvũ) — to decorate}

\label{lesson:GU-C156-sajavvu}



\begin{quote}
Before the new one: say the Gujarati for bindi, then the Gujarati for henna.
\end{quote}

\subsection*{You'll want to know: \gu{સજાવવું}}
\textbf{\gu{સજાવવું}} — \emph{sajāvvũ} — \textquotedblleft{}to decorate\textquotedblright{}.

\begin{quote}
\textbf{\gu{અમે} \gu{દિવાળીમાં} \gu{ઘર} \gu{સજાવીએ} \gu{છીએ}.} — \emph{ame divāḷīmā̃ ghar sajāvīe chīe} — \textquotedblleft{}We decorate the house at Diwali.\textquotedblright{}
\end{quote}

\begin{cousinweb}
The origin of \textbf{\gu{સજાવવું}} (\emph{sajāvvũ}) is not settled. CDIAL lists no Gujarati form of the verb; under Sanskrit \emph{sajja}, \textquotedblleft{}put in place, equipped\textquotedblright{}, it lists Gujarati \emph{sāj}, \textquotedblleft{}dress, equipment\textquotedblright{}, which may be related (CDIAL 13095).
\end{cousinweb}

\begin{grammarlens}[title={What you've built}]
One more word for everyday life at A2.
\end{grammarlens}

\subsection*{Your turn}
\begin{itemize}
\raggedright
  \item \emph{Say it:} \emph{sajāvvũ}
  \item \emph{Say it:} \emph{sajāvvũ}, once more
  \item \emph{Say it:} \emph{ame divāḷīmā̃ ghar sajāvīe chīe} — We decorate the house at Diwali.
  \item \emph{Recall:} say the Gujarati for bindi, then the Gujarati for henna, then say \emph{sajāvvũ} again
\end{itemize}

\begin{tcolorbox}[breakable,colback=teal!4,colframe=teal!35!black,title={Before you move on}]
What does \gu{સજાવવું} mean? (\textquotedblleft{}To decorate\textquotedblright{}.) Say it once more.
\end{tcolorbox}
\section[\texorpdfstring{\gu{બાંધવું} (bā̃dhvũ)}{બાંધવું (bā̃dhvũ)}]{\gu{બાંધવું} (bā̃dhvũ) — to tie, to bind}

\label{lesson:GU-C156-badhvu}



\begin{quote}
Before the new one: say the Gujarati for henna, then the Gujarati for to decorate.
\end{quote}

\subsection*{You'll want to know: \gu{બાંધવું}}
\textbf{\gu{બાંધવું}} — \emph{bā̃dhvũ} — \textquotedblleft{}to tie, to bind\textquotedblright{}.

\begin{quote}
\textbf{\gu{તે} \gu{વાળ} \gu{બાંધે} \gu{છે}.} — \emph{te vāḷ bā̃dhe che} — \textquotedblleft{}She ties up her hair.\textquotedblright{}
\end{quote}

\begin{cousinweb}
\textbf{\gu{બાંધવું}} (\emph{bā̃dhvũ}) is a tadbhava, a word inherited from Sanskrit and reshaped by sound change. It comes from \emph{bandhati}, \textquotedblleft{}binds\textquotedblright{}, and CDIAL glosses the Gujarati verb \textquotedblleft{}to bind, tie\textquotedblright{} (CDIAL 9139).
\end{cousinweb}

\begin{grammarlens}[title={What you've built}]
One more word for everyday life at A2.
\end{grammarlens}

\subsection*{Your turn}
\begin{itemize}
\raggedright
  \item \emph{Say it:} \emph{bā̃dhvũ}
  \item \emph{Say it:} \emph{bā̃dhvũ}, once more
  \item \emph{Say it:} \emph{te vāḷ bā̃dhe che} — She ties up her hair.
  \item \emph{Recall:} say the Gujarati for henna, then the Gujarati for to decorate, then say \emph{bā̃dhvũ} again
\end{itemize}

\begin{tcolorbox}[breakable,colback=teal!4,colframe=teal!35!black,title={Before you move on}]
What does \gu{બાંધવું} mean? (\textquotedblleft{}To tie, to bind\textquotedblright{}.) Say it once more.
\end{tcolorbox}
